Towards Modelling and Reasoning Support forEarly-Phase Requirements EngineeringEric S. K. YuFaculty of Information Studies, University of TorontoToronto, Ontario, Canada M5S 3G6AbstractRequirements are usually understood as stating whata system is supposed to do, as opposed to how it shoulddo it. However, understanding the organizational con-text and rationales (the “Whys”) that lead up to systemsrequirements can be just as important for the ongoing suc-cess of the system. Requirements modelling techniquescan be used to help deal with the knowledge and reason-ing needed in this earlier phase of requirements engineer-ing. However, most existing requirements techniques areintended more for the later phase of requirements engi-neering, which focuses on completeness, consistency, andautomated verification of requirements. In contrast, theearly phase aims to model and analyze stakeholder inter-ests and how they might be addressed, or compromised, byvarious system-and-environment alternatives. This paperargues, therefore, that a different kind of modelling andreasoning support is needed for the early phase. An out-line of the framework is given as an example of a stepin this direction. Meeting scheduling is used as a domainexample.1 IntroductionRequirements engineering (RE) is receiving increasingattention as it is generally acknowledged that the earlystages of the system development life cycle are crucialto the successful development and subsequent deploymentand ongoing evolution of the system. As computer systemsplay increasingly important roles in organizations, it seemsthere is a need to pay more attention to the early stages ofrequirements engineering itself (e.g., [6]).Much of requirements engineering research has takenas starting point the initial requirements statements, whichexpress customer’s wishes about what the system shoulddo. Initial requirements are often ambiguous, incomplete,inconsistent, and usually expressed informally. Many re-quirements languages and frameworks have been proposedfor helping make requirements precise, complete, and con-sistent (e.g., [4] [19] [13] [15]). Modelling techniques(from boxes-and-arrows diagrams to logical formalisms)with varying degrees of analytical support are offered toassist requirements engineers in these tasks. The objec-tive, in these “late-phase” requirements engineering tasks,is to produce a requirements document to pass on (“down-stream”) to the developers, so that the resulting systemwould be adequately specified and constrained, often in acontractual setting.Considerably less attention has been given to support-ing the activities that precede the formulation of the initialrequirements. These “early-phase” requirements engineer-ing activities include those that consider how the intendedsystem would meet organizational goals, why the system isneeded,whatalternativesmightexist, whatthe implicationsof the alternatives are for various stakeholders, and how thestakeholders’ interests and concerns might be addressed.The emphasis here is on understanding the “whys” that un-derlie system requirements [37], rather than on the preciseand detailed specification of “what” the system should do.This earlier phase of the requirements process can bejust as important as that of refining initial requirementsto a requirements specification, at least for the followingreasons:System development involves a great many assump-tions about the embedding environment and task do-main. As discovered in empirical studies (e.g., [11]),poor understanding of the domain is a primary causeof project failure. To have a deep understanding abouta domain, one needs to understand the interests andpriorities and abilities of various actors and players, inaddition to having a goodgrasp of the domain conceptsand facts.Users need help in coming up with initial require-ments in the first place. As technical systems increasein diversity and complexity, the number of technicalalternatives and organizational configurations madepossible by them constitute a vast space of options.A systematic framework is needed to help developersunderstand what users want and to help users under-stand what technical systems can do. Many systemsthat are technically sound have failed to address realneeds (e.g., [21]).Systems personnel are increasingly expected to con-tribute to business process redesign. Instead of au-tomating well-establishedbusiness processes, systemsare now viewed as “enablers” for innovative businesssolutions (e.g., [22]). More than ever before, require-ments engineers need to relate systems to business andorganizational objectives.
Dealing with change is one of the major problemsfacing software engineering today. Having a repre-sentation of organizational issues and rationales inrequirements models would allow software changesto be traced all the way to the originating source –the organizational changes that leads to requirementschanges [18].Having well-organized bodies of organizational andstrategic knowledge would allow such knowledge tobe shared across domains at this high level, deepeningunderstanding about relationships among domains.This would also facilitate the sharing and reuse ofsoftware (and other types of knowledge) across thesedomains.As more and more systems in organizations intercon-nect and interoperate, it is increasingly important tounderstand how systems cooperate (with each otherand with human agents) to contribute to organiza-tional goals. Early phase requirements models thatdeal with organizational goals and stakeholder inter-ests cut across multiple systems and can provide aview of the cooperation among systems within an or-ganizational context.Support for early-phase RE. Early-phase RE activitieshave traditionally been done informally, and without muchtool support. As the complexity of the problem domainincreases, it is evident that tool support will be neededto leverage the efforts of the requirements engineer. Aconsiderable body of knowledge would be built up duringearly-phase RE. This knowledge would be used to sup-porting reasoning about organizational objectives, system-and-environment alternatives, implications for stakehold-ers, etc. It is important to retain and maintain this body ofknowledge in order to guide system development, and todeal with change throughout the system’s life time.A number of frameworks have been proposed to repre-sent knowledge and to support reasoning in requirementsengineering (e.g., [19] [13] [27] [17] [12] [5] [35]). How-ever, these frameworks have not distinguished early-phasefrom late-phase RE. The question then is: Are there mod-elling and reasoning support needs that are especially rel-evant to early-phase RE? If there are specific needs, canthese be met by adapting existing frameworks?Most existing requirements techniques and frameworksare intended more for the later phase of requirements engi-neering, which focuses on completeness, consistency, andautomated verification of requirements. In contrast, theearly phase aims to model and analyze stakeholder inter-ests and how they might beaddressed, or compromised, byvarious system-and-environment alternatives. In this paperit is argued that, because early-phase RE activities have ob-jectives and presuppositions that are different from those ofthe late phase, it would be appropriate to provide differentmodelling and reasoning support for the two phases. Nev-ertheless, a number of recently developed RE techniques,such as agent- and goal-oriented techniques (e.g., [16] [14][17] [12] [7]) are relevant, and may be adapted for early-phase RE.The recently proposed framework [39] is used in thispaper as an example to illustrate the kinds of modelling fea-tures and reasoning capabilities that might be appropriatefor early-phase requirements engineering. It introduces anontology and reasoning support features that are substan-tially different from those intended for late-phase RE (e.g.,as developed in [15]).Section2 reviewsthe framework andoutlines some ofits features, using meeting schedulingas a domainexample.Section 3 discusses, in lightof the experience of developing, the modelling and support requirements for early-phaserequirements engineering. Section 4 reviews related work.Section 5 draws some conclusions from the discussionsandidentifies future work.2 The i* modelling framework for early-phase requirements engineeringThe framework was developed for modelling andreasoning about organizational environments and their in-formation systems [39]. It consists of two main modellingcomponents. TheStrategicDependency(SD)model isusedto describethe dependency relationshipsamong various ac-tors in an organizational context. The Strategic Rationale(SR) model is used to describe stakeholder interests andconcerns, and how they might be addressed by variousconfigurations of systems and environments. The frame-work builds on a knowledge representation approach toinformation system development [27]. This section offersan overview of some of the features of , using primarilya graphical representation. A more formal presentation ofthe framework appears in [39]. The framework has alsobeen applied to business process modelling and redesign[41] and to software process modelling [38].The central concept in is that of the intentional ac-tor [36]. Organizational actors are viewed as having in-tentional properties such as goals, beliefs, abilities, andcommitments. Actors depend on each other for goals tobe achieved, tasks to be performed, and resources to befurnished. By depending on others, an actor may be ableto achieve goals that are difficult or impossible to achieveon its own. On the other hand, an actor becomes vulner-able if the depended-on actors do not deliver. Actors arestrategic in the sense that they are concerned about opportu-nities and vulnerabilities, and seek rearrangements of theirenvironments that would better serve their interests.2.1 Modelling the embedding of systems in orga-nizational environments – the Strategic De-pendency modelConsider a computer-based meeting scheduler for sup-porting the setting up of meetings . The requirementsmight state that for each meeting request, the meetingThe name i* refers to the notion of distributed intentionality whichunderlies the framework.An early version of the framework was presented in [36].The example used in this paper is a simplified version of the oneprovided in [34].
scheduler should try to determine a meeting date and loca-tion so that most of the intended participants will partici-pate effectively. The system would find dates and locationsthat are as convenient as possible. The meeting initiatorwould ask all potential participants for information abouttheir availability to meet during a date range, based on theirpersonal agendas. This includes an exclusion set – dateson which a participant cannot attend the meeting, and apreference set – dates preferred by the participant for themeeting. The meeting scheduler comes up with a proposeddate. The date must not be one of the exclusion dates, andshould ideally belong to as many preference sets as pos-sible. Participants would agree to a meeting date once anacceptable date has been found.Many requirements engineering frameworks and tech-niques have been developed to help refine this kind of re-quirements statements to achieve better precision, com-pleteness, and consistency. However, to develop systemsthat will truly meet the real needs of an organization, oneoften needs to have a deeper understanding of how thesystem is embedded in the organizational environment.For example, the requirements engineer,before proceed-ing to refine the initial requirements, might do well to in-quire:Why is it necessary to schedule meetings ahead oftime?Whydoes themeeting initiator needto askparticipantsfor exclusion dates and preferred dates?Why is a computer-based meeting scheduler desired?And whose interests does it serve?Is confirmation via the computer-based scheduler suf-ficient? If not, why not?Are important participants treated differently? If so,why?Most requirements models are ill-equipped to help an-swer such questions. They tend to focus on the “what”rather than the “why”. Having answers to these “why”questions are important not only to help develop successfulsystems in the first instance, but also to facilitate the de-velopment of cooperation with other systems (e.g., projectmanagement systems and other team coordination “group-ware” for which meeting information may be relevant), aswell as the ongoing evolution of these systems.To providea deeperlevelof understandingabouthowtheproposed meeting scheduler might be embedded in the or-ganizational environment, the Strategic Dependency modelfocuses on the intentional relationships among organiza-tional actors. By noting the dependencies that actors haveon one another, one can obtain a better understanding ofthe “whys”.Considerfirstthe organizationalconfigurationbeforetheproposed system is introduced (Figure 1). The meeting ini-tiator depends on meeting participants to attend meeting. If some participant does not attend the meeting, themeeting initiator may fail to achieve some goal (not madeDMeetingInitiatorExclusionDates(p)PreferredDates(p)ProposedDate(m) Agreement(m,p)  AttendsMeeting(ip,m)Assured(AttendsMeeting(ip,m))  AttendsMeeting(p,m) ImportantParticipantXISADDDDDDDDDDDDD MeetingParticipantDDDDDDDDDepender DependeeLEGENDO  XOpen (uncommitted) CriticalResource DependencyTask DependencyGoal DependencySoftgoal DependencyFigure 1: Strategic Dependency model for meetingscheduling, without computer-based schedulerexplicit in the SD model), or at least not succeed to thedegree desired. This is the reason for wanting to schedulethe meeting in advance. To schedule meetings, the initiatordepends on participants to provide information about theiravailability – in terms of a set of exclusion dates and pre-ferred dates. (For simplicity, we do not separately considertime of day or location.) To arrive at an agreeable date, par-ticipants depend on the initiator for date proposals. Onceproposed, the initiator depends on participants to indicatewhether they agree with the date. For important partici-pants, the meeting initiator depends critically (marked withan “X” in the graphical notation) on their attendance, andthus also on their assurance that they will attend.Dependency types are used to differentiate among thekinds of relationships between depender and dependee, in-volving different types of freedom and constraint. Themeeting initiator’s dependency on participant’s attendanceat the meeting ( ) is a goal depen-dency. It is up to the participant how to attain that goal. Anagreement on a proposed date is mod-elled as a resource dependency. This means that the par-ticipant is expected only to give an agreement. If there isno agreement, it is the initiator who has to find other dates(do problem solving). For an important participant, theinitiator critically depends on that participant’s presence.The initiator wants the latter’s attendance to be assured(). This is modelledas a softgoal dependency. It is up to the depender to decidewhat measures are enoughfor him to be assured,e.g., a tele-phone confirmation. These types of relationships cannot beexpressed or distinguished in non-intentional models thatare used in most existing requirements modelling frame-works.Figure 2 shows an SD model of the meeting schedul-ing setting with a computer-based meeting scheduler. Themeeting initiator delegates much of the work of meetingscheduling to the meeting scheduler. The initiator nolonger needs to be bothered with collecting availability in-formation from participants, or to obtain agreements aboutproposed dates from them. The meeting scheduler alsodetermines what are the acceptable dates, given the avail-
DMeetingInitiator ProposedDate(m) Agreement(m,p)  AttendsMeeting(ip,m)Assured(AttendsMeeting(ip,m))  AttendsMeeting(p,m) MeetingBeScheduled(m) MeetingScheduler ImportantParticipantXISADDDDDDDDDDDDDDD    EnterDateRange     (m)    EnterAvailDates     (m)   MeetingParticipantDDDDDDDDDepender DependeeLEGENDO  XOpen (uncommitted) CriticalResource DependencyTask DependencyGoal DependencySoftgoal DependencyFigure 2: Strategic Dependency model for meetingscheduling with computer-based schedulerability information. The meeting initiator does not carehow the scheduler does this, as longer as the acceptabledates are found. This is reflected in the goal dependency offrom the initiator to the scheduler.The scheduler expects the meeting initiatorto enter the daterange by following a specific procedure. This is modelledvia a task dependency.Note that it is still the meeting initiator who dependson participants to attend the meeting. It is the meetinginitiator (not the meeting scheduler) who has a stake inhaving participants attend the meeting. Assurance fromimportant participants that they will attend the meeting istherefore not delegated to the scheduler, but retained as adependency from meeting initiatorto important participant.TheSD modelmodels themeetingschedulingprocess interms of intentional relationships among agents, instead ofthe flow of entities among activities. This allows analysisof opportunity and vulnerability. For example, the abilityof a computer-based meeting scheduler to achieve the goalof representsan opportunityfor themeeting initiator not to have to achieve this goal himself.On the other hand, the meeting initiator would become vul-nerable to the failure of the meeting scheduler in achievingthis goal.2.2 Modelling stakeholder interests and ratio-nales – the Strategic Rationale modelThe Strategic Dependency model provides one level ofabstraction for describing organizational environments andtheir embedded information systems. It shows external(but nevertheless intentional) relationships among actors,while hiding the intentional constructs within each actor.As illustrated in the preceding section, the SD model canbe useful in helping understand organizational and systemsconfigurations as they exist, or as proposed new configura-tions.During early-phase RE, however, one would also like tohave more explicit representation and reasoning about ac-tors’ interests, and how these interests might be addressedor impacted by different system-and-environment configu-rations – existing or proposed.In the framework, the Strategic Rationale model pro-vides a more detailed level of modelling by looking “in-side” actors to model internal intentional relationships. In-tentional elements (goals, tasks, resources, and softgoals)appear in the SR model not only as external dependencies,but also as internal elements linkedby means-ends relation-ships and task-decompositions (Figure 3). The SR modelin Figure 3 thus elaborates on the relationships between themeeting initiator and meeting participant as depicted in theSD model of Figure 1.For example, for the meeting initiator, an internal goalis that of . This goal can be met(represented via a means-ends link) by scheduling meet-ings in a certain way, consisting of (represented via task-decomposition links): obtaining availability dates frompar-ticipants, finding a suitable date (and time) slot, proposinga meeting date, and obtaining agreement from the partici-pants.These elements of the task are rep-resented as subgoals, subtasks, or resources depending onthe type of freedom of choice as tohow to accomplish them(analogous to the SD model). Thus ,being a subgoal, indicates that it can be achieved in dif-ferent ways. On the other hand, andrefer to specific ways of accomplish-ing these tasks. Similarly, , beingrepresented as a goal, indicates that the meeting initiatorbelieves that there can be more than one way to achieve it(to be discussed in section 2.4, Figure 4).is itself an element of thehigher-level task of organizing a meeting. Other subgoalsunder that task might include equipment be ordered, or thatreminders be sent (not shown). This task has two addi-tional elements which specify that the organizing of meet-ings should be done quickly and not involve inordinateamounts of effort. These qualitative criteria are modelledas softgoals. These would be used to evaluate (and alsoto help identify) alternative means for achieving ends. Inthis example, we note that the existing way of schedulingmeetings is viewed as contributing negatively towards theand softgoals.On the side of the meeting participants, they are ex-pected to do their part in arranging the meeting, and thento attend the meeting. For the participant, arranging themeeting consists primarily of arriving at an agreeable date.This requires them to supply availability information to themeeting initiator, and then to agree to the proposed dates.Participants want selected meeting times to be convenient,and want meeting arranging activities not to present toomany interruptions.The SR model thus provides a way of modelling stake-holder interests, and how they might be met, and the stake-holders evaluation of various alternatives with respect totheir interests. Task-decomposition links provide a hierar-chical description of intentional elements that make up aroutine. The means-ends links in the SR provides under-standing about why an actor would engage in some tasks,
Agreeable(Meeting,Date)AgreeTo DateConvenient(Meeting,Date)ProposedDateAgreementObtainAgreementObtainAvailDates   MergeAvailDates AttendMeetingParticipateInMeetingScheduleMeetingMeetingInitiatorQuality(ProposedDate)UserFriendlyArrangeMeeting MeetingParticipantAttendsMeetingDDDDDDDD+−Quick MeetingBe ScheduledOrganizeMeeting−−+−++PreferredDatesExclusionsDatesDD  MinInterruptionLowEffortLowEffortFindAgreeableDate  FindSuitable  SlotLEGENDGoalTaskResourceSoft−GoalTask−Decompo−   sition linkMeans−ends linkContribution to    softgoalsActorActor Boundary+−Figure 3: A Strategic Rationale model for meeting scheduling, before considering computer-based meeting schedulerpursue a goal, need a resource, or want a softgoal. From thesoftgoals, one can tell why one alternative may be chosenover others. For example, availability information in theform of exclusion sets and preferred sets are collected soas to minimize the number of rounds and thus to minimizeinterruption to participants.2.3 Supporting analysis during early-phase REWhile requirements analysis traditionally aims to iden-tify and eliminate incompleteness, inconsistencies,and am-biguities in requirements specifications, the emphasis inearly-phase RE is instead on helping stakeholders gain bet-ter understanding of the various possibilities for using in-formation systems in their organization, and of the impli-cations of different alternatives. The models offer anumber of levels of analysis, in terms of ability, workabil-ity, viability and believability. These are detailed in [39]and briefly outlined here.When a meeting initiator has a routine to organize ameeting, we say that he is able to organize a meeting. Anactor who is able to organize one type of meeting (say, aproject group meeting) is not necessarily able to organizeanother type of meeting (e.g., the annual general meetingfor the corporation). One needs to know what subtask,subgoals, resources are required, and what softgoals arepertinent.Given a routine, one can analyze it for workability andviability. Organizing meeting is workable if there is aworkable routine for doing so. To determine workability,one needs to look at the workability of each element – forexample, that the meeting initiator can obtaining availabil-ity information from participants, canfind agreeable dates,and can obtain agreements from participants. If the work-ability of an element cannot be judged primitively by theactor, then it needs to be further reduced. If the subgoalis not primitively workable, it willneedto beelaborated intermsofa particularwayfor achiev-ing it. For example, one possible means for achieving itis to do an intersection of the availability information fromall participants. If this task is judged to be workable, thenthe goal node would be workable. Atask can be workable by way of external dependencies onother actors. The workability of andare evaluated in terms of the workabil-ity of the commitment of meeting participants to providesavailability information and agreement. A more detailedcharacterization of these concepts are given in [39].A routine that is workable is not necessarily viable. Al-though computing intersection of time slots by hand is pos-sible, it is slow and error-prone. Potentially good slots maybe missed. When softgoals are not satisficed, the routineis not viable. Note that a routine which is not viable fromone actor’s perspective may be viable from another actor’sperspective. For example, the existing way of arrangingfor meetings may be viable for participants, if the resultingmeeting dates are convenient, and the meeting arrangementefforts do not involve too much interruption of work.The assessment of workability and viability is based onmany beliefs and assumptions. These can be provided asjustifications for the assessment. The believability of therationale network can be analyzed by checking the networkof justifications for the beliefs. For example, the argumentthat “finding agreeable dates by merging available dates”is workable may be justified with the assertion that themeeting initiator has been doing it this way for years, andit works. The belief that meeting participants will supplyavailability information and agree to meeting dates may be
Agreeable(Meeting,Date)AgreeTo DateConvenient(Meeting,Date)ProposedDateAgreementObtainAgreementObtainAvailDates  FindAgreeable  Slot   MergeAvailDatesQuick MeetingBe ScheduledScheduleMeeting AttendMeetingParticipateInMeetingLetScedulerScheduleMeetingScheduleMeetingMeetingInitiatorMeetingSchedulerMeetingBe ScheduledFindAgreeableDateUsingSchedulerFindAgreeableDateByTalkingToInitiatorQuality(ProposedDate)UserFriendlyArrangeMeetingRicherMedium  EnterAvailDatesOrganizeMeeting MeetingParticipantAttendsMeeting    EnterDateRangeDDDDDDDDDDDD+−−++−++−++LEGENDGoalTaskResourceSoft−GoalTask−Decompo−   sition linkMeans−ends linkContribution to    softgoalsActorActor Boundary+−LowEffortLowEffortFigure 4: Strategic Rationale model for a computer-supported meeting scheduling configurationjustified by the belief that it is in their own interests to doso (e.g., programmers who want their code to pass a re-view). The evaluation of these goal graphs (or justificationnetworks) is supported by graph propagation algorithmsfollowing a qualitative reasoning framework [8] [42].2.4 Supporting design during early-phase REDuring early-phase RE, the requirements engineer as-sists stakeholders in identifying system-and-environmentconfigurations that meet their needs. This is a process ofdesign on a higher level than the design of the technicalsystem per se. In analysis, alternatives are evaluated withrespect to goals. In design, goals can be used to help gen-erate potential solutions systematically.In , the SR model allows us to raise ability,workabil-ity, and viability as issues that need to be addressed. Usingmeans-ends reasoning, these issues can be addressed sys-tematically, resulting in new configurations that are then tobe evaluated and compared. Means-ends rules that encodeknowhow in the domain can be used to suggest possible al-ternatives. Issues and stakeholders that are cross-impactedmay be discovered during this process, and can be raisedso that trade-offs can be made. Issues are settled whenthey are deemed to adequately addressed by stakeholders.Once settled, one can then proceed from the descriptivemodel of the framework to a prescriptive model thatwould serve as the requirements specification for systemsdevelopment. Believability can also be raised as an issue,so that assumptions would be justified.In analyzing the SR model of Figure 3, it isfound that the meeting initiator is dissatisfied with theOne approach to this is described in [40].amount of effort needed to schedule a meeting, andhow quickly a meeting can be scheduled. These areraised as the issues and.Sincethe meetinginitiator’s existingroutinefor schedul-ing meetings is deemed unviable, one would need to lookfor new routines. This is done by raising the meeting initia-tor’s ability to schedule meetings as an issue. To addressthis issue, one could try to come up with solutions withoutspecial assistance, or one could look up rules (in a knowl-edge base) that may be applicable. Suppose a rule is foundwhose purpose is and whose howattribute is .This represents knowledge that the initiator has aboutsoftware scheduler systems, their abilities, and their plat-form requirements. The rule helps discover that the meet-ing initiator can delegate the subgoal of meetingschedulingto the (computer-based) meeting scheduler. This consti-tutes a routine for the meeting initiator.Using a meeting scheduler, however, requires partici-
pants to enter availability information in a particular for-mat. This is modelled as a task dependency on participants(an SD link). A routine that provides for this is sought inthe participant. Again, rules may be used to assist in thissearch.When new configurations are proposed, they may bringin additional issues. The new alternatives may have as-sociated softgoals. The discovery of these softgoals canalso be assisted with means-ends rules. For example, us-ing computer-based meeting scheduling may be discov-ered to be negative in terms of medium richness and user-friendliness. These in turn have implications for the effortinvolved for the participant, and the quality of the proposeddates. These newly raised issues also need to be addressed.Once new routines have been identified, they are analyzedfor workability and viability. Further routines are searchedfor until workable and viable ones are found.3 The modelling and reasoning supportneeds of early-phase REIn the preceding section, the framework was outlinedin order to illustrate the kind of modelling and reasoningsupport that would be useful during the early phase of re-quirements engineering. This section summarizes and dis-cusses these modelling and support needs in more generalterms, drawing from the experience of the framework.Knowledge representation and reasoning. Althoughthe example in the preceding section relies primarily oninformal graphical notations, it is clear that a realistically-sized application domain would involve large numbers ofconcepts and relationships. A more formal knowledge rep-resentation scheme would be needed to support modelling,analysis, and design activities. Maintaining a knowledgebase of the knowledge collected and used during early-phase RE is also crucial in order to reap benefits for sup-porting ongoing evolution (e.g., [8]), and for reuse acrossrelated domains.Many of the knowledge-based techniques developed forother phases of software engineering are also applicablehere. For example, knowledge structuring mechanismssuch as classification, generalization, aggregation, and time[20] are equally relevant in early-phase as inlate-phase RE.On the other hand, early-phase RE has certain needs thatare quite distinct from late-phase RE.Degree of formality. While representing knowledge for-mally has the advantage of amenability to computer-basedtool support, the nature of the early-phase suggests thatformality should be used judiciously. The early-phase REprocess is likely to be a highly interactive one, with thestakeholders as the source of information as well as thedecision maker. The requirements engineer acts primarilyin a supporting role. The degree of formality for a sup-port framework therefore needsto reflect this relationship.Use of knowledge representation can facilitate knowledgemanagement and reasoning. However, one should not tryto over formalize, as one may compromise the style ofreasoning needed.One approach is to introduce weaker constructs, suchas softgoals, which requires judgemental inputs from timeto time in the reasoning process, but which can be struc-tured and managed nonetheless within the overall knowl-edge base [7] [39]. The notion of softgoal draws on theconcept of satisficing [33], which refers to finding solu-tions that are “good enough”.Incorporating intentionality. One of the key needs indealing with the subject matter in the early phase seems tobe the incorporation of the concept of the intentional actorinto the ontology. Without intentional concepts such asgoals, one cannot easily deal with the “why” dimension inrequirements.Anumber ofrequirements engineeringframeworkshaveintroduced goal-oriented and agent-oriented techniques(e.g., [16] [14] [17] [12] [7]). In adapting these techniquesfor early-phase RE, one needs to recognize that the focusduring the early phase is on modelling (i.e., describing) theintentionality of the stakeholders and players in the orga-nizational environment. When new alternatives are beingsought (the “design” component in early phase RE), it isthe intentionality of the stakeholders that are being exer-cised. The requirements engineer is helping stakeholdersfind solutions to their problems. The decisions rests withthe stakeholders.In most goal-oriented frameworks in RE, the intention-ality is assumed to be under the control of the requirementsengineer. Therequirements engineermanipulatesthe goals,andmakes decisionson appropriatesolutions tothesegoals.This may be appropriate for late-phase RE, but not for theearly phase.By the end of the early-phase, the stakeholders wouldhave made the major decisions that affect their strategicinterests. Requirements engineers and developers can thenbe given the responsibility to fill in the details and to realizethe system.One consequence of the early/late phase distinction isthat intentionality is harder to extract and incorporate intoa model in the early phase than in the late phase. Stake-holder interests and concerns are typically not readily ac-cessible. The approach adopted in is to introduce thenotion of intentional dependencies to provide a level of ab-straction that hides the internal intentional contents of anactor. The Strategic Dependency model provides a usefulcharacterization of the relationships among actors that is atan intentional level(as opposed to non-intentional activitiesand flows), without requiring the modeller to know muchabout the actors’ internal intentional dispositions. Onlywhen one needs to reason about alternative configurationswould one need to make explicit the goals and criteria forsuch deliberations (in the Strategic Rationale model). Evenhere, the model of internal intentionality is not assumedto be complete. The model typically contains only thoseconcerns that are voiced by the stakeholders in order forthem to achieve the changes they desire.Multi-lateral intentional relationships. In modellingthe embedding of a system in organizational environments,it is necessary to describe dependencies that the systemhas on its environment (human agents and possibly other
systems), as well as the latter’s dependencies on the sys-tem. When the system does not live up to the expectationsof agents in its environment, the latter may fail to achievecertain goals. The reverse can also happen. During early-phase RE, one needs to reason about opportunities andvulnerabilities from both perspectives. Both the systemand its environment are usually open to redesign, withinlimits. When opportunities or vulnerabilities are discov-ered, further changes can be introduced on either side totake advantage of them or to mitigate against them. Amodelling framework for the early-phase thus needs to beable to express multi-lateral intentional relationships and tosupport reasoning about their consequences.In most requirements frameworks, the requirementsmodels are interpreted prescriptively. They state what asystem is supposed to do. This is appropriate for late-phaseRE. Requirements documents are often used in contrac-tual settings – developers are obliged to design the systemsin order to meet the specifications. Once the early-phasedecisions have settled, a conversion from the multi-lateraldependency model to a unilateral prescriptive model for thelate-phase can be made.Distributed intentionality. Another distinctive featureof the early-phase subject matter is that the multiple ac-tors in the domain all have their own intentionality. Ac-tors exercise intentionality (e.g., they pursue goals) in thecourse of their daily routines. Actors have multiple, some-times conflicting, sometimes complementary goals. Theintroduction of a computer system may make certain goalseasier to achieve and others harder to achieve, thus per-turbing the network of strategic dependencies. Differ-ent system-and-environment configurations can thereforebe seen as different ways of re-distributing the pattern ofintentionality . The boundaries may shift (the responsibil-ity for achieving certain goalsmay be delegated from someagents to other agents, some of which may be computersystems), but the actors remain intentional. The processof system-and-environment redesign does not solve all theproblems (i.e., does not (completely) reduce intentional el-ements, such as goals, to non-intentional elements, such asactions). It merely rearranges the terrain inwhich problemsappear and need to be addressed.In contrast, in late-phase RE and in the rest of systemdevelopment, one does attempt to fully reduce goals toimplementable actions.Means-ends reasoning. In order to model and supportreasoning about “why”, and to help come up with alterna-tive solutions, some form of means-ends reasoning wouldappear necessary. However, a relatively weaker form ofreasoning than customarily used in goal-oriented frame-works is needed. This is because of the higher degree ofincompletenessin earlyphase RE.The emphasisis on mod-elling stakeholders’ rationales. Alternative solutions maybe put forth as suggestions, but it is the stakeholders whodecide. The modelling may proceed both “upwards” and“downwards” (from means to ends or vice versa). There isno definitive “top” (since there may always be some highergoal) nor “bottom” (since there is no attempt to purge in-Hence the name i*.tentionality entirely). It is the stakeholders’ decision asto when the issues have been adequately explored and asufficiently satisfactory solution found.The type of reasoning support desired is therefore closerto those developed in issue-based informationsystems, ar-gumentation frameworks, and design rationales (e.g., [10][30] [26] [25]). The approach is an adaptation of aframework developed for dealing with non-functional re-quirements [7], which draws on these earlier frameworks.Organizational actors. In modelling organizational en-vironments, a richer notion of actor is needed. differ-entiates actors into agents, roles, and positions [39]. Inlate-phase requirements engineering, where the focus is onspecifying behaviours rather than intentional relationships,such distinctions may not be as significant. Viewpointshas been recognized as an important topic in requirementsengineering (e.g., [29]). In the early phase, the need totreat multiple viewpoints involving complex relationshipsamong various types of actors is even more important.4 Related workIn the requirements modelling area, the need to modelthe environment is well recognized (e.g., [4] [19] [23] ).Organization and enterprise models have been developedin the areas of organizational computing (e.g., [1]) andenterprise integration (e.g., [9]). However, few of thesemodels have considered the intentional, strategic aspectsof actors. Their focus has primarily been on activitiesand entities rather than on goals and rationales (the “what”rather than the “why”).Anumber ofrequirements engineeringframeworkshaveintroduced concepts of agents or actors, and employ goal-oriented techniques. The framework of [5] uses multiplemodels to model actors, objectives, subject concepts andrequirements separately, and is close in spirit to theframework in many ways. The WinWin framework of [2]identifies stakeholder interests and links them to quality re-quirements. The notion of inquiry cycle in [31] is closelyrelated to the early-phase RE notion, but takes a scenariosapproach. The KAOS framework [12] [35] for require-ments acquisition employs the notions of goals and agents,and provides a methodology for obtaining requirementsspecifications from global organizational goals.However, these frameworks do not distinguish betweenthe needs of early-phase vs. late-phase RE. For example,most of them assume a global perspective on goals, whichare reduced, by requirements engineers, in a primarily top-down fashion, fully to actions. These may be contrastedwith the notion of distributed intentionality in , whereagents are assumed to be strategic, whose intentionality areonly partially revealed, who are concerned about oppor-tunities and vulnerabilities, and who seek to advance orprotect their strategic interests by restructuring intentionalrelationships.
5 ConclusionsUnderstanding “why” has been considered an importantpart of requirements engineering since its early days [32].Frameworks and techniques to explicitly support the mod-elling of and reasoning about agents’ goals and rationaleshave recently been developed in RE. In this paper, it wasargued that making a distinction between early-phase andlate-phase RE could help clarify the ways in which theseconcepts and techniques could be applied to different REactivities.The framework was given as an example in whichagent- and goal-oriented concepts and techniques wereadapted to address some of the special needs of early-phaseRE.The proposal to use a modelling framework tailoredspecifically to early-phase RE and aseparate framework forlate-phase RE implies that a linkage between the two kindsof framework is needed [40]. As with other phases in thesoftware development life cycle, the relationship betweenearly and late phase RE is not strictly sequential or eventemporal. Each phase generates and draw on a certain kindof knowledge, which needs to be maintained throughoutthe life cycle for maximum benefit [24] [20] [28]. The ap-plicationof knowledge-basedtechniques to early-phaseREcould potentially bring about a more systematic approachto this often ad hoc, under-supported phase of system de-velopment.Preliminary assessments of the usefulness of mod-elling in a real setting have been positive [3]. Supportingtools and usage methodologies are being developed in anon-going project [42].AcknowledgmentsThe author gratefully acknowledges the many helpfulsuggestions from anonymous referees, Eric Dubois, BrianNixon, and Lawrence Chung, as well as on-going guid-ance from John Mylopoulos, and financial support fromthe Information Technology Research Centre of Ontario,and the Natural Sciences and Engineering Research Coun-cil of Canada.References[1] A. J. C. Blythe, J. Chudge, J.E. Dobson and M.R. Strens,ORDIT: a new methodology to assist in theprocess of elicit-ing and modelling organizational requirements. Proc. Con-ference on Organizational Computing Systems, MilpitasCA, 1993. pp. 216–227.[2] B. Boehm and H. In, Aids for Identifying Conflicts AmongQuality Requirements, IEEE Software, March 1996.[3] L. Briand, W. Melo, C. Seaman and V. Basili, Character-izing and Assessing a Large-Scale Software MaintenanceOrganization, Proc. 17th Int. Conf. Software Engineering.Seattle, WA. 1995.[4] J. A. Bubenko, Information Modeling in the Context ofSystem Development, Proc. IFIP, 1980, pp. 395-411.[5] J. A. Bubenko, Extending the Scope of Information Mod-eling, Proc. 4th Int. Workshop on the Deductive Approachto Information Systems and Databases, Lloret-CostaBrava,Catalonia, Sept. 20-22, 1993, pp. 73-98.[6] J. A. Bubenko, Challenges in Requirements Engineering,Proc. 2nd IEEE Int. Symposium on Requirements Engineer-ing, York, England, March 1995, pp. 160-162.[7] K. L. Chung, Representing and Using Non-FunctionalRequirements for Information System Development: AProcess-Oriented Approach, Ph.D. Thesis, also Tech. Rpt.DKBS-TR-93-1, Dept. of Comp. Sci., Univ. of Toronto,June 1993.[8] L. Chung, B. Nixon and E. Yu, Using Non-Functional Re-quirements to Systematically Support Change, 2nd IEEEInt. Symp. on Requirements Engineering (RE’95), York,England, March 1995.[9] CIMOSA – Open Systems Architecture for CIM, ESPRITConsortium AMICE, Springer-Verlag 1993.[10] J. Conklin and M. L. Begeman, gIBIS: A Hypertext Toolfor Explanatory Policy Discussions, ACM Transactions onOffice Information Systems, 6(4), 1988, pp. 303-331.[11] B. Curtis, H. Krasner and N. Iscoe, A Field Study of theSoftware Design Process for Large Systems, Communica-tions of the ACM, 31(11), 1988, pp. 1268-1287.[12] A. Dardenne, A. van Lamsweerde and S. Fickas, Goal-Directed Requirements Acquisition, Science of ComputerProgramming,20, 1993, pp. 3-50.[13] E. Dubois, J. Hagelstein, E. Lahou, F. Ponsaert and A. Ri-faut, A Knowledge Representation Language for Require-ments Engineering, Proc. IEEE, 74 (10), Oct. 1986, pp.1431 –1444.[14] E. Dubois, A Logic of Action for Supporting Goal-OrientedElaborations of Requirements, Proc. 5th InternationalWorkshop on SoftwareSpecificationand Design,Pittsburgh,PA, 1989, pp. 160-168.[15] Ph. Du Bois, The Albert II Language – On the Design andthe Use of a Formal Specification Language for Require-ments Analysis, Ph.D. Thesis, Department of ComputerScience, University of Namur, 1995.[16] M. S. Feather, Language Support for the Specification andDevelopment of Composite Systems, ACM Trans. Prog.Lang. and Sys. 9, 2, April 1987, pp. 198-234.[17] S. Fickas and R. Helm, Knowledge Representationand Rea-soning in the Design of Composite Systems, IEEE Trans.Soft. Eng., 18, 6, June 1992, pp. 470-482.[18] O.C.Z. Gotel and A.C.W. Finkelstein, An Analysis of theRequirements Traceability Problem, Proc. IEEE Int. Conf.on Requirements Engineering, Colorado Springs, April1994, pp. 94-101.
[19] S. J. Greenspan, J. Mylopoulos, and A. Borgida, CapturingMore World Knowledge in the Requirements Specification,Proc. Int. Conf. on Software Eng., Tokyo, 1982.[20] S. J. Greenspan, J. Mylopoulos and A. Borgida, On For-mal Requirements Modeling Languages: RML Revisited,(invited plenary talk), Proc. 16th Int. Conf. Software Engi-neering, May 16-21 1994, Sorrento, Italy, pp. 135-147.[21] J. Grudin, Why CSCW Applications Fail: Problems in theDesign and Evaln of Organizational Interfaces, Proc. Con-ference on Computer-Supported Cooperative Work 1988,pp. 85-93.[22] M. Hammer and J. Champy, Reengineering the Corpora-tion: A Manifesto for Business Revolution,HarperBusiness,1993.[23] M. Jackson, System Development, Prentice-Hall, 1983.[24] M. Jarke, J. Mylopoulos, J. W. Schmidt and Y. Vassiliou,DAIDA: An Environment for Evolving Information Sys-tems, ACM Trans. Information Systems, vol. 10, no. 1, Jan1992, pp. 1-50.[25] J. Lee, A Decision Rationale Management System: Cap-turing, Reusing, and Managing the Reasons for Decisions,Ph.D. thesis, MIT, 1992.[26] A. MacLean, R. Young, V. Bellottiand T. Moran, Questions,Options, and Criteria: Elements of Design Space Analysis,Human-Computer Interaction, vol. 6, 1991, pp. 201-250.[27] J. Mylopoulos, A. Borgida, M. Jarke and M. Koubarakis,Te-los: Representing Knowledge about Information Systems,ACM Trans. Info. Sys., 8 (4), 1991.[28] J. Mylopoulos, A. Borgida and E. Yu, Representing Soft-ware Engineering Knowledge, Automated Software Engi-neering, to appear.[29] B. Nuseibeh, J. Kramer and A. Finkelstein, Expressing theRelationships Between Multiples Views in RequirementsSpecification, Proc. 15th Int. Conf. on Software Engineer-ing, Baltimore, 1993, pp. 187-196.[30] C. Potts and G. Bruns, Recording the Reasons for DesignDecisions, Proc. 10th Int. Conf. on Software Engineering ,1988, pp. 418-427.[31] C. Potts, K. Takahashi and A. Anton, Inquiry-Based Re-quirements Analysis, IEEE Software, March 1994, pp. 21-32.[32] D. T. Ross and K. E. Shoman, Structured Analysis for Re-quirements Definition, IEEE Trans. Soft. Eng., Vol. SE-3,No. 1, Jan. 1977.[33] H. A. Simon, The Sciences of the Artificial, 2nd ed., Cam-bridge, MA: The MIT Press, 1981.[34] A. Van Lamsweerde, R. Darimont and Ph. Massonet,The Meeting Scheduler Problem: Preliminary Definition.Copies may be obtained from Prof. Van Lamsweerde,Universite Catholique de Louvain, Unite d’Informatique,Place Sainte-Barbe, 2, B-1348 Louvain-la-Neuve,Belgium.(avl@info.ucl.ac.be)[35] A. Van Lamsweerde, R. Darimont and Ph. Massonet, Goal-Directed Elaboration of Requirements for a Meeting Sched-uler: Problems and Lessons Learnt, Proceedings of 2ndIEEE Int. Symposium on Requirements Engineering, York,England, March 1995, pp. 194-203.[36] E. Yu, Modelling Organizations for Information SystemsRequirementsEngineering, Proceedingsof First IEEE Sym-posium on Requirements Engineering, San Diego, Calif.,January 1993, pp. 34-41.[37] E. Yu and J. Mylopoulos, Understanding Why in Require-ments Engineering – with anExample, Workshop onSystemRequirements: Analysis, Management, and Exploitation,Schloß Dagstuhl, Saarland, Germany, October 4–7, 1994.[38] E. Yu and J. Mylopoulos, Understanding ‘Why’ in SoftwareProcess Modelling, Analysis, and Design, Proc. 16th Int.Conf. on Software Engineering, Sorrento, Italy, May 1994,pp. 159-168.[39] E. Yu, Modelling Strategic Relationships for ProcessReengineering, Ph.D. thesis, also Tech. Report DKBS-TR-94-6, Dept. of Computer Science, University of Toronto,1995.[40] E. Yu, P. Du Bois, E. Dubois and J. Mylopoulos, FromOrganization Models to System Requirements – A ‘Coop-erating Agents’ Approach, Proc. 3rd Int. Conf. on Coop-erative Information Systems (CoopIS-95), Vienna, Austria,May 1995, pp. 194-204.[41] E. Yu and J. Mylopoulos, From E-R to ‘A-R’ – ModellingStrategic Actor Relationships for Business Process Reengi-neering, Int. Journal of Intelligent and Cooperative Infor-mation Systems, vol. 4, no. 2 & 3, 1995, pp. 125-144.[42] E. Yu, J. Mylopoulos and Y. Lesperance, AI Models forBusiness Process Reengineering, IEEE Expert, August1996, pp. 16-23